\scanpage{120}
Noter que les résultats à la Künneth établis ici
sont valables sans hyp $n$ premier à la car., et
ne dépendent pas du th. de dualité. La formule
de Künneth proprement dite aurait d'abord été
établie par Verdier \deleted{par} en utilisant le th. de dualité
globale.

\margin{\uline{variante}}
Noter que la formule de projection n'est plus
exacte si on y remplace $\mathbb R_!f_*$ par $\mathbb Rf_*$, déjà pour
$M_\bullet=A_{U,X}$, \inserted{$L_\bullet=A_X$, $X=U'\subset Y$ imm. ouverte,}
($U$ ouvert de $X$), et les $H^0$, car
cette formule impliquerait que si $U$ est un ouvert
de $Y$, et $F$ est tel que $F\vert(X-f^{-1}(U))=0$, alors
$f_*(F)\vert(Y-U)=0$, ce qui n'est pas vrai si $f$
non propre\ldots\ \deleted{\unsure{Équ}}
Cependant, une démonstration
presque immédiate (valable pour tt topos, pratiquement,
sous seule réserve d'une finitude de dim ch.)
prouve qu'on a un isom
\[
\mathbb Rf_*(L_\bullet)\overset{\mathbb L}{\otimes}M_\bullet-
\mathbb Rf_*(L_\bullet\overset{\mathbb L}{\otimes}M_\bullet)
\]
lorsque $M_\bullet$ est à cohomologie localement
constante, de t.f. (La restriction ``de type fini''
étant d'ailleurs $\unsure{=}$ inutile si $f$ est
qu. cpt qu. sép.) : par dévissage on se ramène en
effet au cas $M=A_Y$ \unsure{où}
\mspage{119}
